%% Chapter 2b — Multi-Agent Coordination and Collective Blindness
%%
%% Source material:
%%   docs/sources/finite-agent-coordination.tex  (primary)
%%
%% This chapter extends the single-receiver theory to ensembles.  The central
%% results are Group Blindness (cellular floors compound multiplicatively into
%% an organ-level floor that is always positive) and Purpose as Fixed-Point
%% (the ensemble's emergent goal is the omega-limit set of its dynamics, not
%% a stored representation in any individual receiver).

\chapter{Multi-Agent Coordination and Collective Blindness}
\label{ch:coordination}

\section{The Multi-Agent S-Functional}
\label{sec:multi-agent-s}

% Ensemble E = {R_1, ..., R_n} of n bounded receivers.
% Joint S-functional:
%   S(E, x; C) = (1/n) sum_{i=1}^n S(R_i, x; C_i)
% where C_i is the projection of C onto receiver i's action space.
%
% Individual floors: beta_i = S_flat(R_i) > 0 for all i (Floor Positivity, Ch. 2).

% TODO

\section{Group Blindness Theorem}
\label{sec:group-blindness}

% Theorem (Group Blindness).
%   For an organ ensemble E of n cells with individual floors beta_i > 0,
%   the composite epistemic floor is:
%       S_flat(E) = prod_{i=1}^n beta_i / (sum_{i=1}^n beta_i)^{n-1}  > 0.
%
% Proof sketch:
%   The composite floor is the fixed point of the catalytic composition rule
%   kappa(gamma_1 diamond gamma_2) = 1 - (1 - kappa_1)(1 - kappa_2).
%   Since each kappa_i < 1, the product is strictly less than 1, and
%   converting back to floor units gives the formula above.
%
% Consequence: Cellular blindness compounds multiplicatively.
%   The organ does NOT magically know what individual cells do not know.
%   An ensemble of blind receivers is still blind — the composite floor never
%   vanishes for any finite n.
%
% Theorem (Asymptotic Floor).
%   As n -> infinity with beta_i = b for all i,
%       S_flat(E) = b^n / (nb)^{n-1} = b / n^{n-1} -> 0
%   but remains strictly positive for all finite n.

% TODO — full proof

\section{Purpose as Fixed-Point}
\label{sec:purpose-fixed-point}

% The "purpose" of a biological ensemble cannot be a stored representation
% (No Template Theorem, Ch. 8).  It is instead defined dynamically.
%
% Definition (Purpose Cell).
%   The purpose cell C* of ensemble E is the attainable cell satisfying
%       Phi_E(C*) = S_flat(E)
%   where Phi_E is the ensemble dynamics map.
%   Equivalently, C* is the omega-limit set of the ensemble trajectory.
%
% Theorem (Purpose Existence, Banach Fixed-Point).
%   If the ensemble dynamics Phi_E: Omega -> Omega is a contraction
%   (||Phi_E(x) - Phi_E(y)|| <= kappa ||x - y|| with kappa < 1), then
%   there exists a unique fixed point C* = lim_{t->infty} Phi_E^t(x_0)
%   for any initial state x_0.
%
% Theorem (Purpose as Attractor).
%   C* is the unique globally stable fixed point.
%   All trajectories converge to C* regardless of initial state.
%
% Interpretation:
%   "Thriving" is not a template that cells possess — it is what the ensemble
%   gravitates toward.  The blind CAN lead the blind to the destination
%   precisely because the destination is not a memorised map but a dynamic
%   equilibrium.

% TODO

\section{Motivation Heterogeneity Theorem}
\label{sec:motivation-heterogeneity}

% Theorem (Motivation Heterogeneity).
%   S_flat(E) depends only on {beta_i}_{i=1}^n, not on the goal contents
%   (G_i, A_i) of individual receivers.
%
% Proof: The composite floor formula (Section~\ref{sec:group-blindness})
%   contains only beta_i.  Goal vectors G_i and action sets A_i enter
%   only through the trajectory Phi_E but not through the floor.
%
% Corollary (No Molecular Specificity Argument).
%   Richer molecular goals (more specific G_i) or more diverse action sets A_i
%   cannot reduce S_flat(E).  The epistemic floor is thermodynamic, not
%   representational.

% TODO

\section{Common-Cell Convergence and Belief Incompatibility}
\label{sec:common-cell}

% Theorem (Common-Cell Convergence).
%   Two agents with representation-disjoint beliefs (incompatible internal
%   models) can nonetheless attain the same action-cell C.
%
% Proof: By Mode Non-Privilege (Ch. 2), any mode achieving S < tau(C)
%   reaches C.  Agents with incompatible beliefs use different modes;
%   both can independently satisfy S < tau(C).
%
% Theorem (Belief Incompatibility).
%   If two agents A_1, A_2 hold incompatible beliefs (their representation
%   cells are disjoint: C_{A_1} cap C_{A_2} = emptyset), their S-values
%   for a shared target C satisfy:
%       S(A_1, x; C) and S(A_2, x; C) are independent.
%
% Theorem (Reachability Monotonicity).
%   If C_1 subseteq C_2, then any ensemble reaching C_1 also reaches C_2.
%
% Theorem (Purpose Stability).
%   C* is Lipschitz-stable under small perturbations of {beta_i}:
%       ||C*(beta') - C*(beta)|| <= L ||beta' - beta||.

% TODO
